\section{Desigualdad de McDiarmid}

¿Como hariamos ahora para estimar funciones de las muestras? 

Por ejemplo, si querriamos:

\[
f(X) = \max_{i} \hat{\mu}_i
\]

¿Que tan bien lo estariamos haciendo?

Observemos que si cambiamos una única observación \(X_j^a\) por otro valor \(X_j^{a\prime}\in[0,1]\).

La media empírica del brazo \(a\) cambia a lo sumo
\[
|\hat\mu(a)-\hat\mu'((a)|
=
\frac{|X_j^a-X_j^{a\prime}|}{n}
\le\frac1n.
\]

Las medias de los demás brazos permanecen iguales. Usando desigualdad triangular obtenemos
\[
|f(x)-f(x')|\le\frac{1}{n}.
\]

Esto motiva el metodo de las diferencias acotadas de McDiarmid, que nos permite acotar la probabilidad de que una función de variables aleatorias independientes se desvíe de su valor esperado.

\begin{theorem}{Desigualdad de McDiarmid (\cite{mcdiarmid1989method})}{teo:mcdiarmid}
	Sea \(X_1, \ldots, X_n\) variables aleatorias independientes y \(f: \mathcal{X}^n \to \mathbbm{R}\) una función tal que para cada \(i \in \{1, \ldots, n\}\) y
	puntos \(x_1, \ldots, x_n, x_i' \in \mathcal{X}\) se cumple que:

	\begin{equation}\label{eq:estabilidad_mcdiarmid}
		\left| f(x_1, \ldots, x_i, \ldots, x_n) - f(x_1, \ldots, x_i', \ldots, x_n) \right| \leq c_i
	\end{equation}

	Luego, si \(f(S) = f(X_1, \ldots, X_n)\), luego \(\forall \varepsilon > 0\):

	\begin{align}
		P\left( f(S) - \E{f(S)} \geq \varepsilon \right)  & \leq \exp\left( -\frac{2\varepsilon^2}{\sum_{i=1}^n c_i^2} \right) \\
		P\left( f(S) - \E{f(S)} \leq -\varepsilon \right) & \leq \exp\left( -\frac{2\varepsilon^2}{\sum_{i=1}^n c_i^2} \right)
	\end{align}

\end{theorem}

\begin{proof}
    Sea \(\mathcal F_k=\sigma(X_1,\ldots,X_k)\), con \(\mathcal F_0\) trivial. Definimos
    \begin{align}
        V_k &= \E{f(S)|\mathcal F_k}-\E{f(S)|\mathcal F_{k-1}},\\
        V &= f(S)-\E{f(S)}=\sum_{k=1}^n V_k.
    \end{align}
    Por la propiedad de la torre,
    \[
        \E{V_k|\mathcal F_{k-1}}=0,
    \]
    así que \((V_k)\) es una secuencia de diferencias de martingala.

    Fijado el pasado, sea
    \[
        g_k(x)=\E{f(S)|X_1,\ldots,X_{k-1},X_k=x}.
    \]
    Por independencia, las muestras futuras tienen la misma distribución para todo \(x\). Aplicando \eqref{eq:estabilidad_mcdiarmid} dentro de la esperanza,
    \[
        |g_k(x)-g_k(x')|\leq c_k.
    \]
    Por tanto, con extremos \(\mathcal F_{k-1}\)-medibles,
    \begin{align*}
        L_k&=\inf_x g_k(x)-\E{f(S)|\mathcal F_{k-1}},\\
        U_k&=\sup_x g_k(x)-\E{f(S)|\mathcal F_{k-1}},
    \end{align*}
    tenemos \(L_k\leq V_k\leq U_k\) y \(U_k-L_k\leq c_k\). Aplicando Azuma-Hoeffding a \(\sum_{k=1}^n V_k\) obtenemos ambas cotas. \qedhere
\end{proof}
